%%%%%%%%%%%%%%%%%%%%%%%%%%%%%%%%%%%%%%%%%%%%%%%%%%%%%%%%%%%
%
%--------------CAPÍTULO 3 - TRABALHOS CORRELATOS----------
%
%
\chapter{Trabalhos Correlatos}
\label{cap:trabalhos-correlatos}

Este capítulo situa a proposta em relação à literatura. A Seção~\ref{sec:vc-pecuaria} traça um panorama do uso de visão computacional na pecuária de precisão a partir de revisões sistemáticas da área. A Seção~\ref{sec:deteccao-confinamento} discute trabalhos de detecção de animais em confinamento, com atenção à capacidade dos modelos de generalizar para condições não vistas no treinamento. A Seção~\ref{sec:deteccao-caprinos} trata dos trabalhos voltados a caprinos. A Seção~\ref{sec:especializacao-contexto} reúne trabalhos de fora da pecuária que lidam com a variação das condições de captura por meio de especialização e de Mistura de Especialistas. Em cada seção, os trabalhos são comparados com a proposta deste trabalho.

\section{Visão Computacional na Pecuária de Precisão}
\label{sec:vc-pecuaria}

O aprendizado profundo tornou-se a abordagem predominante nos problemas de visão computacional aplicados à agropecuária e, na maior parte dos casos analisados, supera as técnicas baseadas em características projetadas à mão \cite{kamilaris2018deep}. Na produção animal, \citeonline{nasirahmadi2017} revisaram o uso de visão de máquina na detecção de comportamentos de bovinos e suínos e apontaram como principal vantagem das câmeras o fato de não serem invasivas, ao contrário dos sensores presos ao corpo dos animais. Os autores também listaram a variação de iluminação, a sobreposição entre animais e a sujeira nas instalações entre as principais fontes de erro desses sistemas, fatores que também estão presentes na baia deste trabalho.

Revisões mais recentes mostram que a pesquisa continua concentrada em poucas espécies. Na revisão sistemática de \citeonline{rohan2023deep_livestock_survey} sobre reconhecimento de comportamento de animais de produção por aprendizado profundo, a maioria dos estudos trata de bovinos e suínos. \citeonline{bhujel2024public_datasets} chegaram a conclusão parecida ao analisar 58 bases de imagens públicas para pecuária de precisão: quase metade é de bovinos, seguidos por suínos e aves. Esses autores apontam a falta de bases anotadas de qualidade como um gargalo da área e observam que as bases existentes têm pouca diversidade de ambientes e muitas vezes não registram as condições de captura. Para quem trabalha com caprinos, portanto, construir uma base própria costuma ser inevitável. A base deste trabalho atende em parte a essas lacunas, pois cada imagem está associada à câmera, à sessão de gravação e ao cenário de captura (perspectiva e iluminação), informação usada tanto no treinamento dos especialistas quanto na divisão entre treino e validação. Ela tem, porém, a limitação de diversidade apontada por \citeonline{bhujel2024public_datasets}: vem de uma única baia e de um único rebanho de oito animais.

\section{Detecção de Animais em Instalações de Confinamento}
\label{sec:deteccao-confinamento}

Suínos em baias coletivas são o caso mais estudado de detecção de múltiplos animais em confinamento, e por isso esses trabalhos dão as evidências mais diretas sobre as dificuldades de um detector nesse ambiente. \citeonline{psota2019multipig} propuseram uma rede totalmente convolucional que detecta cada suíno localizando partes do corpo e as associações entre elas, e publicaram uma base com 2.000 imagens e 24.842 animais anotados, vindos de 17 instalações. Para este trabalho, o resultado mais relevante é a diferença entre os dois cenários de avaliação usados pelos autores. Em ambientes presentes no treinamento, o método atingiu precisão acima de 99\% e revocação acima de 96\%; em ambientes e condições de iluminação não vistos no treinamento, os valores caíram para 91\% de precisão e 67\% de revocação. Métricas próximas do máximo, medidas sobre dados parecidos com os de treino, não garantiram o mesmo desempenho em condições novas.

\citeonline{henrich2026pigdetect} retomaram o problema com uma base de detecção (PigDetect) e outra de rastreamento (PigTrack), formadas por imagens de baias em condições reais, com oclusões e baixa visibilidade. Incluir no treino imagens escolhidas por serem difíceis melhorou a detecção em relação ao uso de imagens amostradas aleatoriamente, e os modelos foram avaliados em baias que não apareciam no treinamento. Os dois trabalhos indicam que a forma de amostrar e dividir as imagens afeta tanto o desempenho obtido quanto a confiabilidade da estimativa desse desempenho.

A variação tratada neste trabalho é de outro tipo. O detector não precisa generalizar para baias ou instalações novas; precisa operar em uma única baia sob quatro cenários de captura conhecidos de antemão (duas perspectivas combinadas com imagens diurnas e noturnas em infravermelho), muito diferentes entre si, mas previsíveis. Por isso, as condições de captura são separadas explicitamente e atribuídas a especialistas distintos. Na avaliação, este trabalho usa validação cruzada por blocos temporais com uma zona de exclusão entre treino e validação (Capítulo~4), o que reduz o vazamento entre quadros vizinhos de um mesmo vídeo. Esse protocolo é menos rigoroso que o teste em instalações não vistas de \citeonline{psota2019multipig} e \citeonline{henrich2026pigdetect}, já que blocos diferentes de uma mesma sessão de gravação podem cair em treino e em validação. Para compensar essa limitação, a arquitetura também é avaliada em vídeos de teste independentes.

\section{Detecção e Monitoramento de Caprinos}
\label{sec:deteccao-caprinos}

Há menos trabalhos com caprinos, e em geral eles tratam a detecção como etapa intermediária de uma tarefa maior. \citeonline{zhang2022goat_counting} propuseram um método de contagem automática de caprinos em baias. Um YOLOv5 modificado, com aumento de dados por RandAugment, uma pirâmide de características adaptativa (AF-FPN) e uma cabeça de detecção dinâmica, detecta as cabeças dos animais; o rastreador DeepSORT acompanha cada detecção ao longo do vídeo e conta os animais que cruzam uma linha virtual. O detector atingiu mAP50 de 92,19\%, contra 84,26\% do YOLOv5 original, e a contagem acertou 98,08\%. As 623 imagens estáticas usadas no treino do detector, porém, foram divididas aleatoriamente entre treino e teste na proporção 7:3, e o artigo não informa se a divisão levou em conta a origem das imagens. Os próprios autores relatam que as detecções perdidas aumentam quando há mais de cerca de 20 animais na área de contagem.

Este trabalho se diferencia de \citeonline{zhang2022goat_counting} em três pontos. O detector localiza o corpo inteiro do animal, e não só a cabeça. A contagem é feita quadro a quadro, pelo número de caixas detectadas, sem rastreamento, de modo que qualquer instabilidade do detector aparece diretamente na contagem. E as imagens são divididas por blocos temporais de cada sessão, e não aleatoriamente, para evitar que imagens quase idênticas apareçam em treino e em teste.

\citeonline{mendez2026goat_behavior} identificaram comportamentos de caprinos (comer, beber, ficar em pé e deitar) em uma instalação experimental de produção intensiva, usando vídeos de uma câmera zenital gravados com e sem luz. Os autores compararam 13 modelos de detecção sobre 10.740 anotações. Os modelos da família YOLO (YOLOv8 e YOLOX) tiveram o melhor desempenho geral, e o YOLOX chegou a mAP50 de 0,957. ``Beber'' foi o comportamento mais difícil de detectar, o que os autores atribuem à semelhança visual com outras classes, ao desbalanceamento entre classes e à distorção da lente; corrigir a distorção melhorou a detecção dessa classe.

Em ambiente e perspectiva, esse é o trabalho mais próximo deste, e seus resultados apoiam a escolha da família YOLO para o domínio. As diferenças estão no escopo. \citeonline{mendez2026goat_behavior} usam uma única perspectiva, e um único modelo absorve as variações de iluminação; aqui, as vistas frontal e superior são combinadas, as imagens diurnas e noturnas em infravermelho são tratadas por especialistas distintos, e essa divisão é comparada com um modelo único. Além disso, a tarefa deles é detectar comportamentos, enquanto este trabalho se limita a detectar o animal como classe única, etapa que antecede a classificação comportamental no projeto do qual faz parte.

\section{Especialização por Contexto e Mistura de Especialistas em Detecção}
\label{sec:especializacao-contexto}

Fora da pecuária, a variação das condições de captura já foi tratada como um problema de especialização. \citeonline{mees2016choosing} propuseram uma Mistura de Especialistas profundos para detecção de pessoas por robôs móveis. Cada especialista processa uma modalidade (aparência em RGB, profundidade ou movimento), e uma rede de roteamento atribui pesos às saídas dos especialistas conforme a cena. O sistema se adaptou a mudanças bruscas de iluminação, como a passagem de ambientes internos para externos, e ao borramento causado pelo movimento da câmera, situações em que um especialista isolado falharia. A ideia é a mesma da proposta deste trabalho: a rede de roteamento usa o contexto da cena para decidir em qual especialista confiar. A diferença está na combinação das saídas. Em \citeonline{mees2016choosing}, elas são ponderadas, como na Equação~\ref{eq:moe-soft}; aqui, o roteamento é rígido (\textit{top-1}, Equação~\ref{eq:moe-hard}), e cada especialista é um detector completo treinado com um subconjunto dos cenários de captura.

\citeonline{wang2019universal} estudaram o problema oposto: um único detector capaz de operar em domínios muito diferentes, de faces e placas de trânsito a imagens médicas, sem saber de antemão a que domínio pertence cada imagem. Os autores compararam um banco de detectores especializados, um por domínio, com um detector universal que compartilha parâmetros e usa atenção por domínio. O detector universal superou o banco de especialistas usando cerca de 1,3 vez os parâmetros de um único detector. A especialização, portanto, não traz ganho garantido e precisa ser comparada na prática com um modelo único treinado sobre todos os domínios. Por esse motivo, este trabalho avalia a arquitetura proposta contra um baseline monolítico treinado sob o mesmo protocolo.

Em escala muito maior, \citeonline{riquelme2021vmoe} mostraram que o roteamento esparso pode ser levado para dentro de redes de visão. No V-MoE, camadas de um \textit{Vision Transformer} são substituídas por conjuntos de especialistas, e cada fragmento da imagem (\textit{patch}) é encaminhado a apenas alguns deles. O modelo igualou o estado da arte em reconhecimento de imagens com até metade do custo de inferência. Ali, especialistas e rede de roteamento são treinados juntos, e a divisão do problema é aprendida, sem um significado definido de antemão. Neste trabalho, a divisão entre especialistas é dada pelos cenários de captura, que são conhecidos e interpretáveis, e a rede de roteamento é treinada à parte, como um classificador de cenário.
